% !TEX root = ../main.tex
\chapter{Thực nghiệm và kết quả}
\label{Chapter4}

Chương này trình bày thiết lập thực nghiệm, kết quả đánh giá CoQLLM so với các baseline,
phân tích ablation và thảo luận. Tất cả số liệu trong chương được báo cáo trên một lần
chạy duy nhất (một seed); khoảng tin cậy chưa được tính -- đây là một hạn chế được thừa
nhận (Mục~\ref{sec:han-che}).

\section{Thiết lập thực nghiệm}
\label{sec:thiet-lap-thuc-nghiem}

\subsection{Tập dữ liệu}
\label{ssec:dataset}

Thực nghiệm được thực hiện trên hai tập dữ liệu chuẩn theo quy trình của
CoLLM~\cite{collm2023}.

\textbf{MovieLens-1M}~\cite{movielens2015} gồm khoảng $1$ triệu đánh giá từ $6{.}040$
người dùng cho $3{.}706$ phim, mỗi đánh giá là số sao từ 1 đến 5. Đánh giá được nhị phân
hóa thành nhãn CTR: $\geq 4$ sao được coi là nhấp (click), còn lại là không nhấp.

\textbf{Amazon-Book}~\cite{cora2024} là một tập dữ liệu thưa hơn nhiều so với MovieLens-1M,
theo đúng phân chia dùng trong CoLLM và CoRA. Tập này có khoảng $23$ nghìn người dùng và
$34$ nghìn sản phẩm sách, với độ thưa và tỉ lệ người dùng/sản phẩm nguội cao hơn hẳn -- đây
chính là điều kiện để kiểm chứng cầu nối CF--LLM trên miền khó, nơi tín hiệu cộng tác yếu và
LLM/Q-Former có nhiều dư địa đóng góp. Nhãn CTR được nhị phân hóa tương tự. Việc bổ sung
Amazon-Book cho phép đánh giá khả năng tổng quát của CoQLLM ngoài miền phim và làm nổi
bật vai trò của mục tiêu bảo toàn thứ tự trong từng người dùng (Mục~\ref{ssec:align-rank}).

\subsection{Phân chia dữ liệu và thiết lập đánh giá}
\label{ssec:phan-chia}

Theo quy trình của CoLLM~\cite{collm2023}, dữ liệu được sắp xếp theo thời gian và phân
chia thành tập huấn luyện, kiểm định và kiểm tra theo tỉ lệ 8:1:1. Để đảm bảo so sánh
công bằng, các siêu tham số được chọn dựa trên tập kiểm định và số liệu cuối cùng được
báo cáo trên tập kiểm tra.

\subsection{Mô hình so sánh}
\label{ssec:baseline}

\textbf{Baseline chính:}
\begin{itemize}
    \item \textbf{CoLLM}~\cite{collm2023}: tích hợp tín hiệu cộng tác vào LLM qua MLP pointwise (mỗi véc-tơ cộng tác thành một token); cùng giao
    thức, cùng tập dữ liệu. Đây là baseline chính để đánh giá đóng góp của cầu nối
    Q-Former. Trên MovieLens-1M, luận văn đối chiếu với hai biến thể CoLLM-MF và CoLLM-DIN.
    \item \textbf{BinLLM}~\cite{binllm2024}: mã hóa nhị phân tín hiệu cộng tác vào LLM -- SOTA gần nhất
    theo quy trình này, có mặt trên cả hai tập dữ liệu.
\end{itemize}

\textbf{Baseline bổ sung:}
\begin{itemize}
    \item \textbf{CoRA}~\cite{cora2024}: đại diện hướng \textit{sinh trọng số} (không soft
    token). Trên Amazon-Book, biến thể CoRA-MF được đưa vào bảng so sánh như một điểm tham
    chiếu cho hướng thiết kế khác với CoQLLM (Mục~\ref{ssec:cora-rel}).
    \item \textbf{MF}~\cite{bpr2009} và \textbf{DIN}: các mô hình lọc cộng tác thuần (không LLM),
    làm mốc sàn để đo phần đóng góp thêm của cầu nối CF--LLM.
    \item \textbf{SellaRec}~\cite{sellarec2024}: học tương phản InfoNCE, backbone
    Qwen2-7B-Base; \textit{lưu ý: CoQLLM dùng Vicuna-7B-v1.5, khác họ backbone, nên so sánh
    trực tiếp về số liệu bị giới hạn -- được thảo luận như baseline không đồng nhất}.
    \item \textbf{ILM}~\cite{ilm2024}: Q-Former cho gợi ý; \textit{lưu ý: ILM không báo
    cáo AUC/uAUC cho bài toán CTR, nên không thể so sánh trực tiếp về số liệu; được
    đề cập trong phần thảo luận}.
\end{itemize}

\textbf{Giao thức lấy số liệu baseline.} Trên MovieLens-1M, số liệu CoLLM/BinLLM lấy theo
quy trình CoLLM-parity. Trên Amazon-Book, số liệu CoLLM-MF/BinLLM/CoRA-MF lấy từ lần chạy
lại (re-run) trong CoRA~\cite{cora2024} để đảm bảo cùng phân chia dữ liệu. Cần lưu ý một
khác biệt quy trình: CoQLLM chọn checkpoint tốt nhất theo \textbf{uAUC}, trong khi
CoLLM/CoRA chọn theo \textbf{AUC}; điều này được ghi chú trực tiếp ở bảng kết quả.

\subsection{Chỉ số đánh giá}
\label{ssec:do-do}

Luận văn dùng hai độ đo chính theo quy trình CoLLM:

\begin{itemize}
    \item \textbf{AUC} (Area Under the ROC Curve): đánh giá khả năng phân biệt tổng thể
    trên toàn tập kiểm tra -- xác suất để điểm dự đoán của mẫu dương cao hơn mẫu âm.

    \item \textbf{uAUC} (user-level AUC): AUC được tính riêng trong phạm vi từng người
    dùng rồi lấy trung bình. uAUC đánh giá chất lượng xếp hạng \textit{trong từng người dùng}
    và bất biến với các phép biến đổi đơn điệu theo từng người dùng.
\end{itemize}

Bổ sung: Accuracy và F1 (độ đo phân loại nhị phân), Precision@K và NDCG@K
(độ đo xếp hạng top-K), và cosine similarity trước/sau căn chỉnh (đo mức độ thu hẹp
khoảng cách ngữ nghĩa CF--LLM).

\subsection{Chi tiết cài đặt}
\label{ssec:cai-dat}

\begin{itemize}
    \item \textbf{LLM backbone}: Vicuna-7B-v1.5~\cite{vicuna2023} (tinh chỉnh chỉ dẫn trên
    nền Llama-2) -- cùng lớp backbone Vicuna instruction-tuned mà CoLLM sử dụng
    (Vicuna-7B-v1.1), nhằm bảo đảm so sánh công bằng; đóng băng hoàn toàn; khuôn mẫu chỉ
    dẫn USER/ASSISTANT.
    \item \textbf{MF}: phân rã ma trận BPR~\cite{bpr2009}, chiều ẩn $d = 256$, được huấn
    luyện và đóng băng trước khi tích hợp.
    \item \textbf{Q-Former}: $Q = 8$ truy vấn, $d_{\text{model}} = 768$, $8$ đầu attention,
    $4$ lớp; nhánh văn bản khởi tạo từ \texttt{bert-base-uncased}.
    \item \textbf{LoRA}~\cite{lora2021}: hạng $r = 8$, $\alpha = 16$, dropout $0{,}05$, áp
    dụng trên các lớp $[q,v]$ của LLM.
    \item \textbf{Kỹ thuật SOTA}: điều kiện hóa truy vấn theo người dùng bật trên cả hai tập dữ liệu;
    hàm mất mát bảo toàn thứ tự ($\lambda = 0{,}2$, kèm lấy mẫu theo nhóm người dùng) bật
    trên Amazon-Book, tắt trên MovieLens-1M.
    \item \textbf{Chọn checkpoint}: theo uAUC trên tập kiểm định (khác CoLLM/CoRA chọn theo
    AUC -- xem ghi chú bảng kết quả).
    \item \textbf{Optimizer}: AdamW; learning rate tùy giai đoạn (Bước 2 Giai đoạn 3 dùng
    $3\times10^{-5}$, giảm theo lịch cosine về $3\times10^{-6}$).
    \item \textbf{Phần cứng}: một GPU NVIDIA A100 80GB.
    \item \textbf{Tham số huấn luyện được}: chỉ LoRA (khoảng vài triệu tham số trên các lớp
    $[q,v]$ của backbone $7$ tỉ tham số), cầu nối Q-Former và các lớp chiếu được cập nhật;
    MF và LLM đóng băng. Không có tham số sinh trọng số. Tổng số tham số huấn luyện được rất
    nhỏ so với $7$ tỉ tham số đóng băng của LLM, phản ánh thiết kế tinh chỉnh hiệu quả tham số.
    \item \textbf{Thời gian huấn luyện}: mỗi epoch mất khoảng $4$--$6$ phút trên phần cứng
    nói trên ($0{,}65$--$0{,}84$ giây/iteration tùy giai đoạn).
\end{itemize}

\section{Kết quả tổng thể}
\label{sec:ket-qua-chinh}

Bảng~\ref{tab:ket-qua-chinh} và Bảng~\ref{tab:ket-qua-book} so sánh CoQLLM với CoLLM,
BinLLM và các baseline trên MovieLens-1M và Amazon-Book theo AUC và uAUC.

\begin{table}[ht]
\caption{Kết quả trên MovieLens-1M (quy trình CoLLM-parity). CoQLLM báo cáo hai cấu hình:
\textit{vanilla} (chỉ soft token) và \textit{+ điều kiện hóa} (cấu hình đề xuất).
CoQLLM chọn checkpoint theo uAUC.}
\label{tab:ket-qua-chinh}
\begin{center}
\begin{tabular}{|L{5.5cm}|c|c|}
\hline
\textbf{Phương pháp} & \textbf{AUC} & \textbf{uAUC} \\
\hline
MF (chỉ cộng tác)~\cite{bpr2009} & 0{,}6482 & 0{,}6361 \\
\hline
DIN (chỉ cộng tác) & 0{,}7166 & 0{,}6459 \\
\hline
CoLLM-MF~\cite{collm2023} & 0{,}7295 & 0{,}6875 \\
\hline
CoLLM-DIN~\cite{collm2023} & 0{,}7243 & 0{,}6897 \\
\hline
BinLLM~\cite{binllm2024} & 0{,}7425 & 0{,}6956 \\
\hline
CoQLLM (vanilla) & 0{,}7335 & 0{,}7009 \\
\hline
\textbf{CoQLLM (+ điều kiện hóa)} & \textbf{0{,}7475} & \textbf{0{,}6968} \\
\hline
\end{tabular}
\end{center}
\end{table}

\begin{table}[ht]
\caption{Kết quả trên Amazon-Book. Số liệu CoLLM-MF/BinLLM/CoRA-MF lấy từ lần chạy lại
trong CoRA~\cite{cora2024} (chọn checkpoint theo AUC); CoQLLM chọn theo uAUC. uAUC tính
trên $2{.}486$ người dùng tập kiểm tra. Run-1 = + điều kiện hóa; Run-2 = thêm
bảo toàn thứ tự.}
\label{tab:ket-qua-book}
\begin{center}
\begin{tabular}{|L{5.5cm}|c|c|}
\hline
\textbf{Phương pháp} & \textbf{AUC} & \textbf{uAUC} \\
\hline
CoLLM-MF~\cite{collm2023} (CoRA re-run) & 0{,}8021 & 0{,}5782 \\
\hline
BinLLM~\cite{binllm2024} (CoRA re-run) & 0{,}8157 & 0{,}5724 \\
\hline
CoRA-MF~\cite{cora2024} (sinh trọng số) & 0{,}8179 & 0{,}6262 \\
\hline
CoQLLM (+ điều kiện hóa, Run-1) & 0{,}8147 & 0{,}5958 \\
\hline
\textbf{CoQLLM (+ bảo toàn thứ tự, Run-2)} & \textbf{0{,}8132} & \textbf{0{,}6062} \\
\hline
\end{tabular}
\end{center}
\end{table}

\textbf{MovieLens-1M.} Cấu hình đề xuất (+ điều kiện hóa) đạt AUC $0{,}7475$ và uAUC
$0{,}6968$, \textit{vượt BinLLM trên cả hai độ đo} ($0{,}7425$/$0{,}6956$) và vượt rõ
CoLLM-MF/CoLLM-DIN. Đây là điểm đáng chú ý nhất: nhiều phương pháp CF--LLM chỉ thắng một
trong hai độ đo. So với cấu hình vanilla, điều kiện hóa truy vấn theo người dùng nâng AUC từ $0{,}7335$
lên $0{,}7475$ ($+0{,}014$) trong khi uAUC gần như không đổi ($0{,}7009 \to 0{,}6968$, chênh
lệch nằm trong biên độ nhiễu) -- đúng như phân tích ở Mục~\ref{sec:nghich-ly}: phần dịch
theo người dùng chủ yếu nâng khả năng phân biệt \textit{giữa} các người dùng (AUC), gần như
trung tính với xếp hạng \textit{trong} người dùng (uAUC).

\textbf{Amazon-Book.} Đây là miền thưa và khó hơn. Với điều kiện hóa truy vấn theo người dùng (Run-1),
CoQLLM đạt AUC $0{,}8147$ và uAUC $0{,}5958$ -- \textit{vượt cả CoLLM-MF và BinLLM (theo lần
chạy lại của CoRA) trên cả hai độ đo}. Bổ sung hàm mất mát bảo toàn thứ tự (Run-2) nâng uAUC
lên $0{,}6062$ ($+0{,}0104$) trong khi AUC gần như đứng yên ($-0{,}0015$) -- đúng thiết kế:
tổn thất này tác động vào thứ tự trong từng người dùng chứ không ảnh hưởng đến khả năng phân biệt
giữa các người dùng.
Con số $0{,}6062$ đưa CoQLLM vào vùng cạnh tranh với CoRA-MF ($0{,}6262$), vốn dùng cơ chế
sinh trọng số khác hẳn.

\textit{Về quy trình so sánh.} CoQLLM chọn checkpoint theo \textbf{uAUC}, còn
CoLLM/CoRA theo \textbf{AUC}; do đó uAUC của CoQLLM là uAUC-tối-ưu-trên-frontier (lợi thế
nhẹ) và cần được đọc kèm khác biệt này. Trên Amazon-Book, số liệu CoLLM/BinLLM có sai khác
đáng kể giữa các lần chạy lại trong tài liệu (uAUC dao động $\approx 0{,}57$--$0{,}63$ tùy
cấu hình huấn luyện và cỡ tập người dùng nhiễu); luận văn vì thế định vị kết quả là ``cạnh
tranh/vượt baseline dưới quy trình có kiểm soát'' thay vì đuổi theo một con số tuyệt đối.

\section{Phân tích khác biệt AUC--uAUC}
\label{sec:nghich-ly}

uAUC tính AUC trong phạm vi từng người dùng rồi lấy trung bình, nên bất biến với các
phép biến đổi đơn điệu riêng theo người dùng. Điều này dẫn đến một đặc điểm quan trọng:
tín hiệu cộng tác nếu chỉ giúp phân biệt giữa các người dùng (thay đổi \textit{mức nền} điểm
dự đoán của mỗi người dùng) thì cải thiện AUC nhưng không cải thiện uAUC. Mô hình cần
học thứ tự xếp hạng \textit{trong} từng người dùng -- cặp sản phẩm nào người dùng này
thích hơn -- thay vì chỉ phân biệt người dùng với nhau.

Đặc điểm này được xác nhận trực tiếp qua phân rã warm/cold (Mục~\ref{ssec:user-group}).
Trên cả hai tập dữ liệu, uAUC ở nhóm người dùng \textit{ấm} (warm, MF có tín hiệu) cao hơn
hẳn nhóm \textit{nguội} (cold): trên Amazon-Book, uAUC warm đạt $0{,}6238$ -- gần chạm
CoRA-MF ($0{,}6262$) -- trong khi uAUC cold chỉ $0{,}5258$, tức xấp xỉ mức ngẫu nhiên. Nhóm
cold kéo trung bình toàn tập xuống $0{,}6062$. Điều này cho thấy trần của uAUC không nằm ở
cầu nối Q-Former (nếu Q-Former yếu thì AUC cũng phải sập, nhưng AUC vẫn cao $\approx 0{,}81$)
mà ở \textit{chất lượng véc-tơ cộng tác của MF cho người dùng/sản phẩm nguội}. Đây là căn cứ để
định hướng cải thiện uAUC về phía MF teacher (Mục~\ref{sec:huong-phat-trien}) thay vì tinh
chỉnh thêm kiến trúc cầu nối.

\section{Thí nghiệm loại bỏ thành phần}
\label{sec:ablation}

Bảng~\ref{tab:ablation} trình bày phân tích ablation các thành phần của CoQLLM.

Bảng~\ref{tab:ablation} tách biệt đóng góp của hai kỹ thuật huấn luyện bằng cách bật/tắt
từng cơ chế trên nền cầu nối soft token, trên cả hai tập dữ liệu.

\begin{table}[ht]
\caption{Ablation hai kỹ thuật huấn luyện. Trên MovieLens-1M: bật điều kiện hóa truy vấn theo người dùng.
Trên Amazon-Book: bật điều kiện hóa truy vấn theo người dùng (Run-1) rồi thêm hàm mất mát bảo toàn thứ tự
(Run-2). Tất cả trên nền soft token; chọn checkpoint theo uAUC.}
\label{tab:ablation}
\begin{center}
\begin{tabular}{|L{2.6cm}|L{4cm}|c|c|}
\hline
\textbf{Tập dữ liệu} & \textbf{Cấu hình} & \textbf{AUC} & \textbf{uAUC} \\
\hline
MovieLens-1M & Soft token (vanilla) & 0{,}7335 & 0{,}7009 \\
\hline
MovieLens-1M & + điều kiện hóa & \textbf{0{,}7475} & 0{,}6968 \\
\hline
Amazon-Book & + điều kiện hóa (Run-1) & 0{,}8147 & 0{,}5958 \\
\hline
Amazon-Book & + bảo toàn thứ tự (Run-2) & 0{,}8132 & \textbf{0{,}6062} \\
\hline
\end{tabular}
\end{center}
\end{table}

\textbf{Đọc bảng ablation.} Hai kỹ thuật tác động lên hai độ đo khác nhau, đúng như thiết kế.
Trên MovieLens-1M, điều kiện hóa truy vấn theo người dùng chủ yếu tác động lên \textit{AUC} (+$0{,}014$), gần trung
tính với uAUC. Trên Amazon-Book, hàm mất mát bảo toàn thứ tự chủ yếu tác động lên \textit{uAUC}
(+$0{,}0104$), gần trung tính với AUC. Sự phân đôi này giải thích vì sao mỗi tập dữ liệu bật
cơ chế khác nhau: trên MovieLens-1M uAUC đã gần bão hòa ($\approx 0{,}70$) nên
hàm mất mát bảo toàn thứ tự không còn dư địa; trên Amazon-Book uAUC còn thấp ($\approx 0{,}59$)
nên tổn thất xếp hạng trong từng người dùng phát huy tác dụng. Việc bật cơ chế theo từng tập dữ
liệu vì thế là một \textit{siêu tham số hợp lệ theo miền}, không phải sự bất nhất về
phương pháp.

\textbf{Các ablation chưa thực hiện.} Do giới hạn tài nguyên, một số ablation vẫn còn để
ngỏ và được ghi nhận minh bạch như hạn chế (Mục~\ref{sec:han-che}): (i)~thay cầu nối Q-Former
bằng một bộ chiếu \textit{attention-MLP} (MLP có self-attention nhưng không có truy vấn học
được) để cô lập đóng góp của \textit{cấu trúc Q-Former} khỏi phần phi tuyến + attention đơn
thuần; (ii)~bật/tắt riêng biểu diễn đa token dưới cùng backbone Vicuna;
(iii)~ablation hàm mất mát bảo toàn thứ tự trên MovieLens-1M (kỳ
vọng trung tính/âm vì uAUC đã bão hòa). Luận văn không điền số ước lượng cho các hàng này.

\section{Phân tích bổ sung}
\label{sec:phan-tich-bo-sung}

\subsection{Khoảng cách ngữ nghĩa}
\label{ssec:cosine-sim}

Để kiểm chứng rằng cầu nối Q-Former thực sự thu hẹp khoảng cách ngữ nghĩa giữa không
gian cộng tác và không gian LLM, luận văn đo cosine similarity giữa:
(i)~véc-tơ cộng tác thô (đầu ra MF) và véc-tơ nhúng token LLM tương ứng trước khi huấn luyện
cầu nối; (ii)~soft token đầu ra Q-Former và véc-tơ nhúng token LLM sau khi huấn luyện.
Kết quả kỳ vọng: cosine similarity tăng sau huấn luyện, cho thấy Q-Former thực sự căn
chỉnh không gian cộng tác về phía không gian ngữ nghĩa LLM.

\subsection{Phân tích theo nhóm warm/cold}
\label{ssec:user-group}

Để kiểm chứng động lực cold-start, luận văn phân rã tập kiểm tra thành hai nhóm: người dùng
\textit{ấm} (warm -- có nhúng MF, tức MF đã ``thấy'' người dùng/sản phẩm trong huấn luyện)
và người dùng \textit{nguội} (cold -- MF không có nhúng, phải dùng đại diện mặc định).
Bảng~\ref{tab:warm-cold} trình bày AUC/uAUC theo từng nhóm trên cả hai tập dữ liệu.

\begin{table}[ht]
\caption{Phân rã warm/cold trên tập kiểm tra (cấu hình đề xuất). Amazon-Book dùng Run-2
(+ bảo toàn thứ tự).}
\label{tab:warm-cold}
\begin{center}
\begin{tabular}{|L{2.6cm}|L{2.4cm}|c|c|c|}
\hline
\textbf{Tập dữ liệu} & \textbf{Nhóm} & \textbf{\#người dùng} & \textbf{AUC} & \textbf{uAUC} \\
\hline
MovieLens-1M & toàn bộ & $7{.}331$ & 0{,}7475 & 0{,}6968 \\
\hline
MovieLens-1M & warm & $4{.}153$ & \textbf{0{,}7787} & \textbf{0{,}7254} \\
\hline
MovieLens-1M & cold & $3{.}178$ & 0{,}7030 & 0{,}5969 \\
\hline
Amazon-Book & toàn bộ & $2{.}486$ & 0{,}8132 & 0{,}6062 \\
\hline
Amazon-Book & warm & $1{.}648$ & \textbf{0{,}8190} & \textbf{0{,}6238} \\
\hline
Amazon-Book & cold & $630$ & 0{,}7928 & 0{,}5258 \\
\hline
\end{tabular}
\end{center}
\end{table}

Kết quả nhất quán trên cả hai tập dữ liệu: nhóm warm đạt uAUC cao hơn hẳn nhóm cold
($0{,}7254$ vs $0{,}5969$ trên MovieLens-1M; $0{,}6238$ vs $0{,}5258$ trên Amazon-Book). Trên
Amazon-Book, uAUC nhóm warm ($0{,}6238$) gần chạm CoRA-MF ($0{,}6262$), trong khi nhóm cold
xấp xỉ mức ngẫu nhiên. Điều này khẳng định uAUC bị chặn trên bởi \textit{chất lượng véc-tơ cộng tác
cho người dùng/sản phẩm nguội} chứ không phải bởi cầu nối: ở nhóm warm -- nơi MF cung cấp tín
hiệu tốt -- CoQLLM khai thác hiệu quả và tiệm cận baseline mạnh nhất. Nhận định này khớp với
phân tích khác biệt AUC/uAUC (Mục~\ref{sec:nghich-ly}) và chỉ ra hướng cải thiện tiếp theo
nằm ở MF teacher trên nhóm cold (Mục~\ref{sec:huong-phat-trien}).

\section{Thảo luận}
\label{sec:thao-luan}

\subsection{Đối chiếu với ILM}
\label{ssec:doi-chieu-ilm}

ILM~\cite{ilm2024} không báo cáo AUC/uAUC cho bài toán CTR, nên không thể so sánh
số liệu trực tiếp. Đây là lý do chính luận văn đặt ra mục tiêu đánh giá Q-Former dưới
quy trình CoLLM: cung cấp cơ sở so sánh công bằng trong cùng điều kiện thực nghiệm.

\subsection{Ưu điểm và hạn chế của Q-Former}
\label{ssec:diem-manh-han-che}

So với mô-đun chiếu MLP của CoLLM, cầu nối Q-Former có cơ chế chọn lọc đặc trưng rõ
ràng hơn qua tập truy vấn học được và self-attention. Vì tín hiệu cộng tác chỉ đi qua đường soft token, LLM
được giữ nguyên khi suy luận -- không có chi phí sinh trọng số theo từng mẫu như hướng sinh
trọng số của CoRA. Đánh đổi là chi phí huấn luyện cầu nối cao hơn lớp MLP đơn giản, cần cân
nhắc tùy điều kiện triển khai.

\section{Hạn chế nghiên cứu}
\label{sec:han-che}

Luận văn thừa nhận các hạn chế thực nghiệm sau:
\begin{enumerate}
    \item \textbf{Một seed:} Số liệu báo cáo từ một lần chạy duy nhất; chưa có khoảng tin
    cậy, kết quả có thể dao động giữa các lần chạy.
    \item \textbf{Thiếu ablation attention-MLP:} Chưa có hàng ablation thay cầu nối
    Q-Former bằng bộ chiếu attention-MLP (không truy vấn học được). Đây là đối chứng quan
    trọng nhất để cô lập đóng góp của \textit{cấu trúc Q-Former} khỏi phần phi tuyến +
    attention đơn thuần; hiện được ghi nhận là hướng thực nghiệm còn để ngỏ.
    \item \textbf{Các ablation thành phần chưa đầy đủ:} Bật/tắt riêng biểu diễn đa token
    và hàm mất mát bảo toàn thứ tự trên MovieLens-1M đều chưa được chạy
    dưới cùng backbone Vicuna; kết luận về từng thành phần vì thế còn hạn chế.
    \item \textbf{Baseline không đồng nhất:} SellaRec dùng backbone Qwen2-Base khác họ với
    Vicuna; so sánh với SellaRec phản ánh khác biệt cả về cơ chế lẫn backbone LLM. Trên
    Amazon-Book, số liệu CoLLM/BinLLM/CoRA lấy từ lần chạy lại của CoRA (chọn theo AUC),
    khác quy trình chọn checkpoint theo uAUC của CoQLLM.
    \item \textbf{Trần uAUC ở nhóm cold:} uAUC bị chặn trên bởi chất lượng nhúng MF cho
    người dùng/sản phẩm nguội (Mục~\ref{ssec:user-group}), không phải bởi cầu nối; cải
    thiện tiếp theo cần MF teacher mạnh hơn trên nhóm cold.
\end{enumerate}
